\chapter{Iteration \& Recursion}
\label{ch:iteration}

In imperative programming paradigms, repetitive operations are typically handled using loop constructs like \texttt{for} or \texttt{while}. These constructs act upon the machine state, explicitly commanding the computer to increment a counter or re-evaluate a condition until a threshold is met. 

However, functional programming forbids the mutation of state. If variables cannot change, how can we increment a loop counter? The answer, deeply rooted in mathematical logic and set theory, is \textbf{Recursion}.

\section{The Functional Loop: Recursion}
Recursion occurs when a function is defined in terms of itself. It is the programmatic equivalent of a mathematical \textit{recurrence relation}. 

Consider the mathematical definition of summing all integers up to $n$, defined as $S(n)$:
\begin{align*}
S(0) &= 0 \quad \text{(Base Case)} \\
S(n) &= n + S(n-1) \quad \text{(Recursive Step)}
\end{align*}

This elegant mathematical definition translates directly into functional code without the need for arbitrary mutating counter variables. Let us examine how this is implemented in Python:

\lstinputlisting[language=Python, caption={Recursion and Lazy Iteration in Python}, label={lst:chapter6_recursion_py}]{../code/python/chapter6/iteration_recursion.py}

Notice how \texttt{sum\_recursive} is a nearly literal translation of the mathematical recurrence relation. The function continues to add frames to the ``call stack'' until it hits the base case ($n=0$), at which point it collapses back, summing the results.

\section{The Pythonic Compromise: Recursion Limits and \texttt{itertools}}
While elegant, deep recursion in Python uncovers a structural limitation: Python imposes a strict recursion depth limit to prevent stack overflow errors. If you attempt to recursively sum to 2000, Python will throw a \texttt{RecursionError}.

This limitation exists because Python lacks \textbf{Tail Call Optimization (TCO)}. 

Therefore, pure functional programmers in Python often lean on an incredibly powerful part of the standard library: \texttt{itertools}. This module provides a suite of tools for processing collections of data using \textbf{Iterators} and \textbf{Generators}. 

Generators calculate values \textit{lazily}—they only compute the next value in the sequence when explicitly asked for it. As shown in Listing \ref{lst:chapter6_recursion_py}, we can define an infinite sequence using \texttt{itertools.count} or a custom generator function yielding Fibonacci numbers indefinitely, but our program won't crash because it pulls from this infinite well only as needed.

\section{Haskell: Native TCO and Infinite Lists}
Unlike Python, strictly pure functional languages like Haskell are built from the ground up to support unbounded recursion.

\lstinputlisting[language=Haskell, caption={Recursion and Infinite Lists in Haskell}, label={lst:chapter6_recursion_hs}]{../code/haskel/chapter6/iteration_recursion.hs}

In Listing \ref{lst:chapter6_recursion_hs}, we see Haskell's pattern matching handle the base case gracefully: \texttt{sumRecursive 0 = 0}. We also introduce the concept of the \textbf{Tail Call}. In \texttt{sumTailRecursive}, the recursive call to the \texttt{go} helper function is the absolute final operation evaluated in that scope. Because there's no lingering work to do after the call, Haskell optimizes the call stack memory away, mutating an accumulator variable silently at the machine-code level while preserving the pure functional interface for the programmer.

Finally, Haskell's extreme dedication to lazy evaluation allows the straightforward declaration of infinite lists: \texttt{[1..]}. The list exists as an unevaluated promise (a ``thunk'') until a function like \texttt{take} forces its evaluation. This allows functional algorithms to describe massive, even endless data structures conceptually without memory explosion.

\section{Tail Recursion: Euclid's GCD vs. Factorial}
\label{sec:tail_recursion}

In our exploration of recursion, we have seen that recursive functions can be extremely elegant translations of mathematical recurrence relations. However, we also encountered a physical constraint in Python: the call stack. When a function calls itself, the runtime system must allocate memory (a stack frame) to keep track of the local variables and the execution state of the calling function.

But not all recursive calls are created equal. Depending on where the recursive call occurs inside the function, compilers and interpreters can perform a powerful optimization called \textbf{Tail Call Optimization (TCO)}. 

\subsection{Defining the Tail Position}

A function call is said to be in the \textbf{tail position} if it is the absolute final action performed by the function before returning. If a function's recursive call is in the tail position, the function is \textbf{tail-recursive}.

Why does this matter? If the recursive call is the last thing a function does, the runtime doesn't need to keep the current stack frame alive. There are no local variables to inspect and no operations left to perform after the sub-call returns. The runtime can simply reuse the current stack frame for the next recursive call. This turns a potentially stack-overflowing recursion into a memory-efficient loop that runs in $O(1)$ space.

Let us compare two classic mathematical algorithms to see the difference between a tail call and a non-tail call.

\subsection{Greatest Common Divisor (Euclid's Algorithm)}

Euclid's algorithm for finding the greatest common divisor (GCD) of two integers $a$ and $b$ is defined recursively as:
\[
\gcd(a, b) = 
\begin{cases} 
a & \text{if } b = 0 \\ 
\gcd(b, a \bmod b) & \text{otherwise} 
\end{cases}
\]

Let's look at the structure of the recursive step: $\gcd(b, a \bmod b)$. When the function makes this recursive call, it does not perform any operation on the result. The value returned by the recursive call is immediately returned to the caller. Therefore, Euclid's algorithm is \textbf{naturally tail-recursive}.

\subsection{Factorial Calculation}

Now, consider the mathematical definition of the factorial of a non-negative integer $n$, denoted as $n!$:
\[
n! = 
\begin{cases} 
1 & \text{if } n = 0 \\ 
n \times (n - 1)! & \text{otherwise} 
\end{cases}
\]

If we translate this recurrence relation directly into a recursive function, the recursive step evaluates $n \times (n-1)!$. 
Here, the recursive call is $(n-1)!$. After this call returns a value, the function must multiply that value by $n$. Because of this pending multiplication operation, the recursive call is \textbf{not} in the tail position. The runtime must retain every single stack frame to remember the multiplier $n$ for each step of the recursion.

\subsection{The Accumulator Pattern}

To make the factorial function tail-recursive, we must eliminate the pending multiplication after the recursive call. We can do this using the \textbf{Accumulator Pattern}.

Instead of waiting for the recursive call to return before multiplying, we multiply first and pass the running product down to the next call using an extra parameter: the **accumulator** ($acc$). 
Mathematically, we can reformulate the factorial relation with an auxiliary function $F(n, acc)$:
\begin{align*}
F(0, acc) &= acc \\
F(n, acc) &= F(n - 1, n \times acc)
\end{align*}
We compute $n!$ by initializing the accumulator to 1: $n! = F(n, 1)$.

Because the recursive call $F(n - 1, n \times acc)$ has no pending operations, it is now in the tail position.

\subsection{Code Implementations}

Let us examine these concepts side-by-side. First, we look at the Python implementations:

\lstinputlisting[language=Python, caption={Tail Recursion and Euclid's GCD in Python}, label={lst:tail_recursion_py}]{../code/python/chapter6/tail_recursion.py}

Although Python does not support TCO automatically (due to design philosophies surrounding stack traces and debugging), writing code in a tail-recursive style is still an important exercise for understanding functional structure. 

In GHC (the Haskell compiler), tail recursion is fully optimized. Let's inspect the Haskell equivalent:

\lstinputlisting[language=Haskell, caption={Tail Recursion and Euclid's GCD in Haskell}, label={lst:tail_recursion_hs}]{../code/haskel/chapter6/tail_recursion.hs}

In Listing \ref{lst:tail_recursion_hs}, we implement the accumulator pattern for the factorial function using a local helper function named \texttt{go}. GHC will recognize both \texttt{factorialTailRecursive} and \texttt{gcdEuclid} as tail-recursive and optimize them into flat loops at compile time, completely eliminating the risk of stack overflows.

